% Ohm's law on a single edge and Kirchhoff's voltage law (KVL) around a triangle loop.
\begin{tikzpicture}[x=1mm,y=1mm,font=\figurefont\fontsize{8}{10}\selectfont,>=Stealth,
  vertex/.style={circle,draw=FigureBlue,fill=FigureBlueFill,inner sep=0pt,
    minimum size=5.2mm,line width=.55pt},
  rvertex/.style={circle,draw=FigureGreen,fill=FigureGreenFill,inner sep=0pt,
    minimum size=5.2mm,line width=.55pt},
  edge/.style={draw=FigureStroke,line width=.6pt},
  flow/.style={->,draw=FigureBlue,line width=.82pt},
  looparrow/.style={->,draw=FigureRed,dashed,line width=.68pt},
  note/.style={font=\FigureNoteFont\fontsize{7}{8.5}\selectfont,text=FigureInk,inner sep=1pt},
  panel/.style={font=\figurefont\fontsize{7.5}{9}\selectfont,text=FigureInk}]
  \definecolor{FigureBlue}{HTML}{7188A3}
  \definecolor{FigureBlueFill}{HTML}{E5EAF0}
  \definecolor{FigureGreen}{HTML}{829C8B}
  \definecolor{FigureGreenFill}{HTML}{E8EEE9}
  \definecolor{FigureYellow}{HTML}{B4A16F}
  \definecolor{FigureYellowFill}{HTML}{F1ECDD}
  \definecolor{FigureRed}{HTML}{A7787D}
  \definecolor{FigureInk}{HTML}{3B4252}
  \definecolor{FigureMuted}{HTML}{6B7280}
  \definecolor{FigureStroke}{HTML}{4C4D4F}

  % Panel (a): Ohm's law on single edge i -> j
  \begin{scope}[xshift=0mm]
    \node[vertex] (i) at (0,10) {$i$};
    \node[rvertex] (j) at (34,10) {$j$};
    \node[note,above=1.5mm] at (i.north) {$v_i=3$};
    \node[note,above=1.5mm] at (j.north) {$v_j=1$};
    \draw[flow] (i) -- node[above,note,text=FigureBlue] {$f_e=c_e(v_i-v_j)=1$}
                       node[below,note,text=FigureMuted] {$\Delta v_e=2,\ c_e=\tfrac{1}{2}$} (j);
    \node[panel] at (17,-7) {(a) 单边 Ohm 定律};
  \end{scope}

  % Panel (b): KVL around loop 1 -> 2 -> 3 -> 1
  \begin{scope}[xshift=56mm]
    \node[vertex] (v1) at (0,18) {$1$};
    \node[vertex] (v2) at (26,18) {$2$};
    \node[rvertex] (v3) at (26,0) {$3$};
    \node[note,left=1.2mm] at (v1.west) {$v_1=\tfrac{1}{2}$};
    \node[note,right=1.2mm] at (v2.east) {$v_2=\tfrac{1}{4}$};
    \node[note,right=1.2mm] at (v3.east) {$v_3=0$};

    \draw[flow] (v1) -- node[above,note] {$\Delta v_{12}=+\tfrac{1}{4}$} (v2);
    \draw[flow] (v2) -- node[right,note] {$\Delta v_{23}=+\tfrac{1}{4}$} (v3);
    \draw[->,draw=FigureRed,line width=.8pt] (v3) -- node[below left,note,text=FigureRed] {$\Delta v_{31}=-\tfrac{1}{2}$} (v1);

    \draw[looparrow] (16,13.5) arc[start angle=60,end angle=-230,radius=4.2mm];
    \node[note,text=FigureRed] at (17.5,11) {$\sum\Delta v=0$};
    \node[panel] at (13,-7) {(b) 闭环 Kirchhoff 电压定律};
  \end{scope}
\end{tikzpicture}
